\section{Introduction}
UV unwrapping projects the 3D surface of a mesh onto a 2D plane, assigning every 3D surface point a corresponding 2D UV coordinate. This step is fundamental in 3D‑content‑creation pipelines because it enables detailed surface information-such as material properties (e.g., base‑color, roughness, normal maps), along with auxiliary maps like ambient occlusion and displacement—to be efficiently stored and manipulated in 2D space.

A principal component of UV unwrapping is surface parameterization, which flattens the 3D surface while trying to preserve geometric properties such as angles and areas. For meshes with complex geometry, flattening the entire surface into a single 2D chart introduces large distortion. Consequently, chart segmentation (or seam cutting) is typically used to divide the mesh into multiple charts along strategically placed seams, allowing each chart to be flattened with reduced distortion. Finally, UV packing arranges the resulting charts within the unit square (the UV atlas) to maximize texture‑space use.

Although UV unwrapping has been studied extensively, existing methods are typically tuned for well-behaved meshes, such as those created by professional 3D artists. They often fail on more complex, AI‑generated meshes. Such meshes, typically extracted from neural‑field isosurfaces (e.g., via marching cubes~\cite{lorensen1998marching}), tend to have bumpy surfaces, many small triangles, and poor geometric quality (e.g., disconnected components or holes). On such data, existing methods may time‑out or return extremely fragmented atlases in which a single chart holds only one or a handful of triangles. This extreme fragmentation hampers texture painting and editing, introduces texture‑bleed and baking or rendering artifacts at chart boundaries, and burdens downstream applications with an unwieldy number of charts.

Some recent approaches mitigate fragmentation~\cite{srinivasan2024nuvo,zhang2024flatten,zhao2025flexpara} by representing the UV mapping using a neural field and optimizing such a field for each 3D shape from scratch. While these methods can effectively control the number of charts generated, they typically run for more than thirty minutes and still exhibit noticeable distortion. Other methods~\cite{li2018optcuts, autoCuts} jointly optimize seam length and distortion but are likewise computationally expensive. Moreover, some existing approaches like~\cite{sorkine2002bounded,lscm, zhou2004iso} segment charts or identify seams using heuristics based on local geometric properties, rather than leveraging the concept of geometric or semantic parts. This can lead to unintuitive or suboptimal chart boundaries that split semantically coherent regions across multiple charts, further complicating downstream tasks such as texture authoring, part-based editing, and semantic-aware rendering.

In this paper, we introduce \ourmodel, a part-based UV unwrapping pipeline for 3D meshes that generates UV mappings with a small number of part-aligned charts while maintaining low distortion, as well as robust and efficient processing—typically completed within a few to several tens of seconds. \ourmodel builds on a recent learning‑based method, PartField~\cite{liu2025partfield}, which produces a hierarchical part tree for the input mesh. \ourmodel also proposes several novel geometric heuristics that further decompose simple local parts into charts that can be flattened with minimal distortion. Combining high‑level semantic decomposition from PartField with fine‑grained geometric cuts, \ourmodel employs a top‑down recursive search that minimizes chart count while keeping each chart’s distortion below a user‑specified threshold. \ourmodel uses established surface parameterization algorithms (e.g., ABF++~\cite{sheffer2005abf++}) for chart flattening and proven packing algorithms for optimal atlas layout. To ensure high speed and robustness, the pipeline incorporates extensive parallelization and acceleration strategies, as well as dedicated handling of non-manifold and degenerate cases.

By explicitly incorporating semantic priors, our approach yields several key benefits. First, semantics improve decomposition by reducing excessive reliance on local geometry, which often causes over-segmentation and long runtimes. Second, semantic cues preserve the coherence of object parts, preventing chart boundaries from cutting through meaningful regions (e.g., across the flat surface of a TV screen, as in Figure~\ref{fig:partuv_logo}, or a human face), thereby facilitating editing and rendering tasks. Third, semantic grouping naturally supports better chart packing strategies, allowing related charts to be organized together—either within a single atlas (top right of Figure~\ref{fig:teaser}) or across multiple atlases (bottom right of Figure~\ref{fig:teaser}). This enhances organizational clarity and simplifies the process of locating and editing related charts. Finally, seams guided by semantic boundaries tend to fall in perceptually unobtrusive locations, improving the visual quality of textured models.

We evaluate \ourmodel on four datasets spanning man-made, AI-generated, CAD models, and common 3D shapes (e.g., Stanford Bunny, XYZ Dragon), and compare it against widely used UV unwrapping tools such as xatlas~\cite{xatlas}, Blender~\cite{blender}, and Open3D~\cite{zhou2018open3d}, as well as the recent neural-based methods~\cite{srinivasan2024nuvo}. As shown in Figure~\ref{fig:teaser}, our method decomposes input meshes into significantly fewer charts—also resulting in shorter seam lengths—while maintaining low angular and area distortion comparable to baseline methods. The incorporation of explicit part priors not only helps the segmented charts better align with part boundaries but also enables new applications, such as packing semantic-aligned parts into separate atlases. Moreover, \ourmodel maintains a high success rate, handles a wide range of meshes (including large, complex, and non-manifold ones), and processes each mesh quickly—typically within a few to several tens of seconds. 


\begin{figure*}
    \centering
    \includegraphics[width=\linewidth]{figures/partuv_pipeline.png}
    \vspace{-2.5em}
    \caption{\textbf{Pipeline of PartUV.} Given a mesh $\mathcal{M}$, we first leverage the learning-based method PartField to predict a part-aware feature field. By applying a clustering algorithm to this field, we obtain a hierarchical part tree $\mathcal{T}$ for the input mesh $\mathcal{M}$. We then recursively traverse the part tree $\mathcal{T}$ starting from the root node. For each visited node $\mathcal{P}$, we apply two novel geometry-based strategies to segment the corresponding part mesh into multiple sets of charts $\mathscr{C}$. Each chart in the set is then flattened using the ABF algorithm~\cite{sheffer2005abf++}, and its distortion is evaluated. If the distortion exceeds a user-specified threshold $\tau$, we recursively traverse the left and right children of the part tree $\mathcal{T}$. Otherwise, we adopt the segmented charts and their corresponding UV mappings for the part mesh.}
    \vspace{-3.5ex}
    \label{fig:pipeline}
\end{figure*}